\section{March 17}

\paragraph{Recall:}
For a line bundle $L$ on a scheme $X$, we have defined:
\[ c_1(L) \cap : \CH_i(X) \to \CH_{i-1}(X), \quad [V] \mapsto c_1(L) \cap [V] \coloneqq \text{ the class of } L|_V. \]

\begin{proposition}\label{prop:properties_of_intersection_with_first_chern_class_operator}
  We have the following properties of the first Chern class operator.
  \begin{enumerate}
    \item Let $L,L'$ be line bundles on a scheme $X$, $\alpha, \alpha' \in \CH_i(X)$.
      Then
      \begin{align*}
        c_1(L \otimes L') \cap \alpha & = c_1(L) \cap \alpha + c_1(L') \cap \alpha, \\
        c_1(L) \cap (\alpha + \alpha') & = c_1(L) \cap \alpha + c_1(L) \cap \alpha', \\
        c_1(L^\vee) \cap \alpha & = - c_1(L) \cap \alpha \in \CH_{i-1}(X).
      \end{align*}
    \item Let $f: X \to Y$ be a proper morphism, $L$ be a line bundle on $Y$ and $\alpha \in \CH_i(X)$.
      Then
      \[ f_*( c_1(f^*L) \cap \alpha ) = c_1(L) \cap f_*\alpha \in \CH_{i-1}(Y). \]
    \item Let $f: X \to Y$ be a flat morphism of relative dimension $n$, $L$ be a line bundle on $Y$ and $\alpha \in \CH_i(Y)$.
      Then
      \[ f^*( c_1(L) \cap \alpha ) = c_1(f^*L) \cap f^*\alpha \in \CH_{i+n-1}(X). \]
    \item Let $L,L'$ be line bundles on a scheme $X$ and $\alpha \in \CH_i(X)$.
      Then
      \[ c_1(L) \cap ( c_1(L') \cap \alpha ) = c_1(L') \cap ( c_1(L) \cap \alpha ) \in \CH_{i-2}(X). \]
  \end{enumerate}
  % \Yang{}
\end{proposition}
\begin{proof}
  The assertions (a)-(c) follow from the corresponding properties of pseudo-divisors in \cref{prop:basic_properties_of_intersection_with_pseudo-divisors}.
  For (d), we can assume $\alpha = [X]$ and $X$ is a variety.
  Then it follows from \cref{thm:commutativity_of_intersection_with_cartier_divisors}.
\end{proof}


\subsection{Vector bundles and projective bundles}

  Let $X$ be a scheme and $\pi:E \to X$ be a (geometric) vector bundle of rank $r$.
  The sheaf $\calE$ of sections of $E$ is a locally free sheaf of rank $r$ on $X$.
  We have $E = \calSpec_X(\Sym^\bullet \calE^\vee)$.
  Define the \emph{projective bundle} $\bbP(E)$ of $E$ by 
  \[ \bbP(E) \coloneqq \calProj_X(\Sym^\bullet \calE^\vee) \]
  with the projection $p:\bbP(E) \to X$.
  The fiber of $p$ over a point $x \in X$ is the projective space $\bbP(E_x)$ parametrizing lines in the vector space $E_x$.

  We have the Euler sequence on $\bbP(E)$:
  \[ 0 \to \calO_{\bbP(E)}(-1) \to p^*E \to \calT_{\bbP(E)/X}(-1) \to 0, \]
  where $\calO_{\bbP(E)}(-1)$ is the tautological line bundle on $\bbP(E)$ and $\calT_{\bbP(E)/X}$ is the relative tangent bundle of $\bbP(E)$ over $X$.

  \begin{remark}
    In \cite{Har77}, he defines the projective bundle $\bbP(\calE) = \calProj_X(\Sym^\bullet \calE)$, which parametrizes hyperplanes in the fibers of $E$.
    Its $\bbP(\calE)$ is the dual projective bundle of our $\bbP(E)$.
  \end{remark}

  \begin{remark}
    The projection $\pi$ and $p$ are flat morphisms of relative dimension $r$ and $r-1$ respectively.
    And $p$ is also proper.
  \end{remark}

  \begin{theorem}\label{thm:Chow_group_of_vector_and_projective_bundle}
    Let $X$ be a scheme, $\pi: E \to X$ be a (geometric) vector bundle of rank $r$ and $p: \bbP(E) \to X$ be the projective bundle of $E$.
    \begin{enumerate}
      \item The pullback $\pi^*: \CH_i(X) \to \CH_{i+r}(E)$ is an isomorphism for all $i$.
      \item We have isomorphisms of abelian groups:
        \[\theta_E : \bigoplus_{j=0}^{r-1} \CH_{i-(r-1)+j}(X) \to \CH_i(\bbP(E)), \quad (a_0, \ldots, a_{r-1}) \mapsto \sum_{j=0}^{r-1} c_1(\calO_{\bbP(E)}(1))^j \cap p^*a_j. \]
        % \Yang{check the grading.}
    \end{enumerate}
  \end{theorem}

  \begin{example}
    Apply \cref{thm:Chow_group_of_vector_and_projective_bundle} to $X = \Spec k$ and $E = \bbA^r$, we get 
    \[ \CH_i(\bbA^r) = \begin{cases}
        \bbZ \cdot [\bbA^r] & i = r, \\
        0 & i \neq r,
    \end{cases} \quad \CH_i(\bbP^{r-1}) = \begin{cases}
        \bbZ \cdot [L_i] & 0 \leq i \leq r-1, \\
        0 & \text{otherwise}.
    \end{cases} \]
    Here $L_i$ is a linear subspace of $\bbP^{r-1}$ of dimension $i$.
  \end{example}

  \begin{example}
    Apply \cref{thm:Chow_group_of_vector_and_projective_bundle} to a smooth curve $X = C$ and $E = C \times \bbA^2$ so that $\bbP(E) = C \times \bbP^1$, we get
    \[ \CH_1(C \times \bbP^1) = p^*\CH_0(C) \oplus c_1(\calO_{\bbP(E)}(1)) \cap p^*\CH_1(C) \]
    and 
    \[ \CH_0(C \times \bbP^1) = c_1(\calO_{\bbP(E)}(1)) \cap p^*\CH_0(C) \xrightarrow{\sim} \CH_0(C). \]

    In general, $\CH_0$ is unchanged by passing to a projective bundle.
    In fact, $\CH_0$ is a stable birational invariant for smooth projective varieties: it is unchanged under birational transformations and stable under taking products with projective spaces.
  \end{example}

  \begin{proof}[Proof of \cref{thm:Chow_group_of_vector_and_projective_bundle}]
    \hspace{1em}
    \begin{step}
      Surjectivity of $\pi^*$.
    \end{step}
    Let $\pi: E \to X$ be as in the statement.
    WTS: $\pi^*$ is an isomorphism.
    Let $U \subseteq X$ be any affine open subset such that $E|_U \cong U \times \bbA^r$.
    Let $Z = X \setminus U$.
    By the localization sequence, we have a commutative diagram with exact rows:
    \[\begin{tikzcd}
        \CH_{i}(Z) \arrow[r]\arrow[d] & \CH_{i}(X) \arrow[r]\arrow[d] & \CH_{i}(U) \arrow[r]\arrow[d] & 0 \\
        \CH_{i+r}(E|_Z) \arrow[r] & \CH_{i+r}(E) \arrow[r]  & \CH_{i+r}(E|_U) \arrow[r]  & 0.
    \end{tikzcd}\]
    By diagram chasing, it's enough to show the assertion for $U$ and $Z$.
    By induction on the dimension of $X$, we can assume the assertion holds for $Z$.
    By induction on $r$, we can assume $r = 1$.
    By induction on the number of irreducible components, we can assume $X$ is irreducible.
    In this case, WTS: $\pi^*: \CH_i(X) \to \CH_{i+1}(X \times \bbA^1)$ is surjective.
    Let $V \subseteq X \times \bbA^1$ be a subvariety of dimension $i+1$.
    WTS, $[V] \in \Image(\pi^*)$.

    % \Yang{dimension of $V$}

    By replacing $X$ by $\overline{\pi(V)}$, we can assume $V$ dominates $X$.
    Hence $V$ defines a closed point on $\Spec \fkk(X) \times_\kk \bbA^1_\kk = \Spec \fkk(X)[t]$.
    Let $\frakp \subset \fkk(X)[t]$ be the corresponding prime ideal.
    And $f \in \fkk(X)[t]$ be a generator of $\frakp$.
    Regarding $f$ as a rational function on $X \times \bbA^1$,
    \[ \divisor(f) = V + \sum_i m_i [D_i], \]
    where $D_i$ are some prime divisors on $X \times \bbA^1$ which do not dominate $X$.
    Hence they are of the form $D_i = D_i' \times \bbA^1$ for some subvariety $D_i' \subseteq X$.
    Hence 
    \[ [V] \sim - \sum n_j \pi^*[D_j'] \in \CH_{i+1}(X \times \bbA^1). \]
    This proves the surjectivity of $\pi^*$.

    \begin{step}
      Surjectivity of the map $\theta_E: \bigoplus_{j=0}^{r-1} \CH_{i-(r-1)+j}(X) \to \CH_i(\bbP(E))$.
    \end{step}
    By localization and noetherian induction, we can assume $E$ is trivial, i.e. $E = X \times \bbA^r$ and $\bbP(E) = X \times \bbP^{r-1}$.
    ETS: letting $F = X \times \bbA^{r+1}$, if $\theta_E$ is surjective, then $\theta_F$ is also surjective.
    We have a diagram:
    \[ \begin{tikzcd}
      \bbP(E) = X \times \bbP^{r-1} \arrow[r, "f", hook] \arrow[dr, "p"'] &  \bbP(F) = X \times \bbP^r \arrow[d, "q"] & \arrow[l, "g"', hook'] E = X \times \bbA^r \arrow[dl, "\pi"] \\
      & X & 
    \end{tikzcd}. \]
    This diagram induces localization sequences:
    \[\begin{tikzcd}
        \CH_i(\bbP(E)) \arrow[r, "f_*"] & \CH_i(\bbP(F)) \arrow[r, "g^*"] & \CH_i(E) \arrow[r] & 0. \\
        & \CH_{i-r}(X) \arrow[u, "q^*"] \arrow[ur, "\pi^*"'] & 
    \end{tikzcd}\]

    Let $\beta \in \CH_i(\bbP(F))$.
    Since $\pi^*$ is surjective, there exists $\alpha \in \CH_{i-r}(X)$ such that $g^*\beta = \pi^*\alpha$.
    By the commutativity of the diagram, $\beta - q^*\alpha \in \Image(f_*)$.
    By induction hypothesis, there exists $(\alpha_j) \in \bigoplus_{j=0}^{r-1} \CH_{i+j-(r-1)}(X)$ such that 
    \[ f_*\theta_E((\alpha_j)) = f_*\left( \sum_{j=0}^{r-1} c_1(\calO_{\bbP(E)}(1))^j \cap p^*\alpha_j \right) = \beta - q^*\alpha. \]
    Note that $\calO_{\bbP(E)}(1) = f^*\calO_{\bbP(F)}(1)$ and $f_*(p^*\alpha_j) = q^*\alpha_j \cap c_1(\calO_{\bbP(F)}(1))$, we have
    \[ f_*\theta_E((\alpha_j)) = \sum_{j=0}^{r-1} c_1(\calO_{\bbP(F)}(1))^j \cap f_*(p^*\alpha_j) = \sum_{j=0}^{r-1} c_1(\calO_{\bbP(F)}(1))^{j+1} \cap q^*\alpha_j. \]
    We are done.

    \begin{step}\label{step:injectivity_of_theta_E_in_the_proof_of_chow_group_of_projective_bundle}
      Injectivity of $\theta_E$.
    \end{step}
    WTS: if $\sum_{j=0}^{r-1} c_1(\calO_{\bbP(E)}(1))^j \cap p^*\alpha_j = 0$, then $\alpha_j = 0$ for all $j$.
    I claim that for any $\alpha \in \CH_i(X)$, we have
    \[ p_*(c_1(\calO_{\bbP(E)}(1))^j \cap p^*\alpha) = \begin{cases}
        0 & j < r-1, \\
        \alpha & j = r-1.
    \end{cases} \]
    Assuming this claim, we have 
    \[ \alpha_j = p_*\left( c_1(\calO_{\bbP(E)}(1))^{r-1-j} \cap \theta_E((\alpha_j)) \right) = 0 \]
    for all $j$.

    For the proof of the claim, we can assume $\alpha = [V]$ and $X = V$ is a variety.
    If $j< r-1$, then $p_*(c_1(\calO_{\bbP(E)}(1))^j \cap p^*[X]) \in \CH_{dim X + j - (r-1)}(X) = 0$.
    If $j = r-1$, then $p_*(c_1(\calO_{\bbP(E)}(1))^{r-1} \cap p^*[X]) = m [X]$. 
    ETS $m = 1$.
    By replacing $X$ by an open subset, we can assume $E$ is trivial.
    Then $c_1(\calO_{\bbP(E)}(1))^{r-1} \cap p^*[X] = \{*\} \times X$ and $p_*(\{*\} \times X) = [X]$.
    Since flat pullback and proper pushforward are commutative in fiber squares, we have $p_*(c_1(\calO_{\bbP(E)}(1))^{r-1} \cap p^*[X]) = [X]$ in general as well.
    % \Yang{}

    \begin{step}
      Injectivity of $\pi^*$.
    \end{step}

    Assuming $\pi^*\alpha = 0 \in \CH_{i+r}(E)$ for some $\alpha \in \CH_i(X)$, we want to show $\alpha = 0$.
    Let $F = E \oplus (\bbA^1 \times X)$.
    We have the following diagram:
    \[\begin{tikzcd}
      \CH_i(\bbP(E)) \arrow[r, "f_*"] & \CH_i(\bbP(F)) \arrow[r, "g^*"] & \CH_i(E) \arrow[r] & 0. \\
      & \CH_{i-r}(X) \arrow[u, "q^*"] \arrow[ur, "\pi^*"'] & 
    \end{tikzcd}\]
    Hence we have $q^*\alpha = f_*\theta_E((\alpha_j))$ for some $(\alpha_j) \in \bigoplus_{j=0}^{r-1} \CH_{i+j-(r-1)}(X)$.
    Expanding $f_*\theta_E((\alpha_j))$, we have
    \begin{align*}
      q^*\alpha & = f_*\left( \sum_{j=0}^{r-1} c_1(\calO_{\bbP(E)}(1))^j \cap p^*\alpha_j \right) \\
      & = \sum_{j=0}^{r-1} c_1(\calO_{\bbP(F)}(1))^j \cap f_*(p^*\alpha_j) \\
      & = \sum_{j=0}^{r-1} c_1(\calO_{\bbP(F)}(1))^{j+1} \cap q^*\alpha_j.
    \end{align*}
    Here $f_*(p^*\alpha_j) = q^*\alpha_j \cap c_1(\calO_{\bbP(F)}(1))$ can be verified by restricting to open dense subset of components of $\alpha_j$ and using the case of trivial bundles.
    Then we have 
    \[ \theta_F((-\alpha, \alpha_0, \ldots, \alpha_{r-1})) = -q^*\alpha + \sum_{j=0}^{r-1} c_1(\calO_{\bbP(F)}(1))^{j+1} \cap q^*\alpha_j = 0. \]
    By the injectivity of $\theta_F$, we have $\alpha = 0$.
  \end{proof}


\subsection{Segre classes of vector bundles}

  We'll first define the ``inverse Chern classes'', called the Segre classes.

  \begin{definition}
    Let $X$ be a scheme, $E$ be a vector bundle of rank $r$ on $X$, $p: \bbP(E) \to X$ be the projective bundle of $E$.
    Define the \emph{$j$-th Segre class homomorphism}  of $E$ by
    \[ s_j(E) \cap : \CH_i(X) \to \CH_{i-j}(X), \quad \alpha \mapsto p_*\left(c_1(\calO_{\bbP(E)}(1))^{r-1+j} \cap p^*\alpha\right). \]
  \end{definition}

  \begin{proposition}
    \begin{enumerate}
      \item Let $X$ be a scheme, $E$ be a vector bundle of rank $r$ on $X$, $\alpha \in \CH_i(X)$.
      Then 
      \[ s_j(E) \cap \alpha = \begin{cases}
          0 & j < 0, \\
          \alpha & j = 0.
      \end{cases} \]
      % \item Let $E,E'$ be vector bundles of rank $r,r'$ on $X$ respectively.
      % Then 
      % \[ s_j(E) \cap ( s_{j'}(E') \cap \alpha ) = s_{j'}( E') \cap ( s_j(E) \cap \alpha ) \]
      % for all $j,j' \in \bbZ$ and $\alpha \in \CH_i(X)$.

      \item Let $f: X \to Y$ be a proper morphism, $E$ be a vector bundle of rank $r$ on $Y$ and $\alpha \in \CH_i(X)$.
      Then
      \[ f_*( s_j(f^*E) \cap \alpha ) = s_j(E) \cap f_*\alpha. \]

      \item Let $f: X \to Y$ be a flat morphism, $E$ be a vector bundle of rank $r$ on $Y$ and $\alpha \in \CH_i(Y)$.
      Then
      \[ f^*( s_j(E) \cap \alpha ) = s_j(f^*E) \cap f^*\alpha. \]

      \item Let $E,E'$ be vector bundles of rank $r,r'$ on $X$ respectively, $\alpha \in \CH_i(X)$.
      Then
      \[ s_j(E) \cap ( s_{j'}(E') \cap \alpha ) = s_{j'}( E') \cap ( s_j(E) \cap \alpha ). \]

      \item Let $L$ be a line bundle on $X$ and $\alpha \in \CH_i(X)$.
      Then $s_1(L) \cap \alpha = - c_1(L) \cap \alpha$.
    \end{enumerate}
  \end{proposition}
  \begin{proof}
    (a) is the claim in \cref{step:injectivity_of_theta_E_in_the_proof_of_chow_group_of_projective_bundle} in the proof of \cref{thm:Chow_group_of_vector_and_projective_bundle}.
    % \Yang{}
    
    (c) - (d) follow from the corresponding statements for $c_1(L) \cap$ for line bundles $L$ on $X$.
    For instance, as for (b), we have a commutative diagram:
    \[\begin{tikzcd}
      \bbP(f^*E) \arrow[r, "f'"] \arrow[d, "p'"] & X \arrow[d, "f"] \\
      \bbP(E) \arrow[r, "f''"] & Y
    \end{tikzcd}\]
    Then 
    \begin{align*}
      f_*( s_j(f^*E) \cap \alpha ) & = f_*( c_1(\calO_{\bbP(f^*E)}(1))^{r-1+j} \cap p'^*\alpha ) \\
      & = c_1(\calO_{\bbP(E)}(1))^{r-1+j} \cap f''_*p'^*\alpha \\
      & = c_1(\calO_{\bbP(E)}(1))^{r-1+j} \cap q^*f_*\alpha \\
      & = s_j(E) \cap f_*\alpha.
    \end{align*}
    % \[ f_*( s_j(f^*E) \cap \alpha ) = f_*( c_1(\calO_{\bbP(f^*E)}(1))^{r-1+j} \cap p'^*\alpha ) = c_1(\calO_{\bbP(E)}(1))^{r-1+j} \cap f''_*p'^*\alpha = s_j(E) \cap f_*\alpha. \]
    % \Yang{}
    For (e), we have $\bbP(L) = X$ and $\calO_{\bbP(L)}(-1) = L$.
    Hence $s_1(L) \cap \alpha = c_1(\calO_{\bbP(L)}(1)) \cap p^*\alpha = -c_1(L) \cap \alpha$.
  \end{proof}
